\section*{ЗАКЛЮЧЕНИЕ}
\addcontentsline{toc}{section}{ЗАКЛЮЧЕНИЕ}

В ходе выполнения данной работы было разработано веб-приложение для управления авторским контентом с интеграцией интеллектуального ассистента. Система представляет собой веб-приложение на фреймворке Django с использованием базы данных SQLite, текстового редактора Quill и интеграцией с API OpenRouter.

Разработанное приложение позволяет авторам структурировать текст, редактировать с автосохранением, импортировать черновики из PDF/DOCX, экспортировать главы в DOCX и получать помощь от ИИ-ассистента.

Основные результаты работы:

\begin{enumerate}
    \item Проведён анализ предметной области, выявлены проблемы традиционного подхода к организации писательской работы, обоснована необходимость автоматизации. Выполнен обзор существующих систем, на основе которого сделан вывод об отсутствии полноценного веб-решения с интеграцией ИИ-ассистента.
    
    \item Сформулировано техническое задание: определены функциональные требования в виде вариантов использования, а также нефункциональные требования. Разработаны макеты пользовательского интерфейса.
    
    \item Спроектирована клиент-серверная архитектура, обоснован выбор технологий. Разработана реляционная база данных, обеспечивающая хранение иерархической структуры. Спроектированы маршруты API и пользовательский интерфейс.
    
    \item Реализованы все программные компоненты системы, разработаны ключевые алгоритмы, а также проведено функциональное тестирование, подтвердившее корректную работу всех сценариев использования.
\end{enumerate}

Все требования, объявленные в техническом задании, были полностью реализованы, все задачи, поставленные в начале разработки проекта, решены. Разработанное веб-приложение может быть использовано авторами для написания и структурирования книг, импорта черновиков, экспорта готовых глав и получения помощи от ИИ-ассистента.
